\documentclass[12pt]{letter}
\usepackage[utf8]{inputenc}
\usepackage{geometry}
\usepackage{hyperref}
\hypersetup{hidelinks}
\geometry{margin=2.5cm}
\signature{Shuaidong Gao\\Chongqing Institute of Foreign Studies\\Qijiang Campus, Chongqing 401420, China\\gsd3247186514@gmail.com\\ORCID: 0009-0004-5641-3581}

\begin{document}

\begin{letter}{Editor-in-Chief\\\textit{Bioinformatics}\\Oxford University Press}

\opening{Dear Editor,}

I am writing to submit ``CAGate: Cluster-Aware Gating for Differentiable Causal Discovery in Cancer Genomics'' for consideration in \textit{Bioinformatics}.

The idea came from a simple observation: methods like NOTEARS are effective at learning causal graphs from data, but on real cancer transcriptomics they collapse once you go beyond about 150 genes. That threshold cuts out most of the interesting biology. The cluster structure in tumor data---molecular subtypes, immune infiltration states, copy-number-driven expression programs---is well known, yet nobody seemed to be asking whether it could stabilize causal discovery in the high-dimensional regime. So I tried building clustering directly into the NOTEARS optimization loop. The results were stronger than I expected.

CAGate runs adaptive clustering on model residuals at each outer iteration, then computes a MAD-aware gate weight per cluster ($\gamma_k = 1 - \exp(-\alpha \cdot \overline{\mathrm{MAD}}_k)$). Homogeneous subpopulations drive edge discovery; heterogeneous ones get down-weighted. Crucially, clustering and causal discovery co-evolve---there is no separate cluster-then-discover pipeline. This feedback loop is what makes the difference.

Across all 33 TCGA cancer types (68 configurations spanning three dimensions and two methods), CAGate consistently finds more causal edges than NOTEARS: mean $\Delta = +158$ edges, with the largest gains on small-sample cancers ($\Delta = +213$ for $n < 323$; Mann--Whitney $U = 222$, $p = 0.0006$). At $d = 200$, CAGate recovers $552 \pm 34$ edges where NOTEARS finds zero. The most striking result---one I only fully appreciated after benchmarking every type---is that the method works best on the smallest datasets, precisely where you would expect it to fail. External validation against CTD (89--98\%), ClinGen (69--82\%), and STRING confirms the recovered edges are enriched for real biology, not noise.

I believe \textit{Bioinformatics} is the right home for this work. It is fundamentally a methods paper with large-scale biological validation: a lightweight extension to a known framework that opens up a previously inaccessible data regime. The method adds $<2\%$ runtime overhead and the default hyperparameters worked on every cancer without adjustment. Complete code, TCGA download scripts, pre-computed results, and 10 figure-generation pipelines are at \url{https://github.com/gsd3247186514-cell/CAGate-Replication}; every figure can be reproduced with a single command.

I should mention that I am an independent researcher---this work was done solo over the past year at Chongqing Institute of Foreign Studies, without funding or collaborators. The CAGate method and its extensions are the subject of dual patent applications with the China National Intellectual Property Administration and the United States Patent and Trademark Office (both filed 2026). The manuscript has not been submitted elsewhere, and I have no competing interests to declare.

I suggest the following five reviewers, all with complementary expertise and no institutional ties to the NOTEARS baseline group:
\begin{enumerate}
\item Caroline Uhler (MIT, \href{mailto:cuhler@mit.edu}{cuhler@mit.edu}) --- causal discovery in genomics
\item Sach Mukherjee (Cambridge/DZNE, \href{mailto:sach.mukherjee@mrc-bsu.cam.ac.uk}{sach.mukherjee@mrc-bsu.cam.ac.uk}) --- causal structure learning for biomedicine
\item Casey S.\ Greene (U.\ Colorado, \href{mailto:casey.s.greene@cuanschutz.edu}{casey.s.greene@cuanschutz.edu}) --- gene regulatory network inference
\item Trey Ideker (UC San Diego, \href{mailto:tideker@ucsd.edu}{tideker@ucsd.edu}) --- cancer systems biology
\item Jonas Peters (ETH Zurich, \href{mailto:jonas.peters@stat.math.ethz.ch}{jonas.peters@stat.math.ethz.ch}) --- causal inference methodology
\end{enumerate}

Thank you for your time and consideration.

\closing{Sincerely,}

\end{letter}
\end{document}
